% Source: OpenAI2026AreaLaw, 06-transport.tex, lines 28--34 and 99--128.
% Manuscript revision: adc7f1241b42e322a6451854ab7e4b4c146bf78a.
\section{Intertwining a weighted mean tree}
\label{sec:mean_tree_intertwine}

Let $T$ be a finite weighted binary tree as in
Definition~\ref{def:mean_tree}, with leaf labels in an arbitrary set $I$.
For a family of positive definite matrices $(A_i)_{i\in I}$, denote its
root by $M_T(A)$. Thus a leaf labelled $i$ has root $A_i$, and a vertex
with parameter $p\in[0,1]$ and subtrees $l,r$ has root
\begin{align}
  M_T(A)&=F_p\bigl(M_l(A),M_r(A)\bigr),
\end{align}
where $F_p$ is the binary geometric-mean formula of
Theorem~\ref{thm:geometric_mean_value_intertwine}.
The following identity is auxiliary to the weighted-tree construction
and Lemma~7.1 of~\cite{OpenAI2026AreaLaw}.

\begin{theorem}[Intertwining the root of a mean tree]
  \label{thm:mean_tree_intertwine}
  \lean{Matrix.MeanTree.eval_intertwine}
  \leanok
  \uses{def:mean_tree}
  Let $(A_i)_{i\in I}$ and $(B_i)_{i\in I}$ be families of positive
  definite matrices on two finite-dimensional complex Hilbert spaces.
  Suppose that a linear map $J$ from the second space to the first satisfies,
  for every $i\in I$,
  \begin{align}
    A_iJ&=JB_i.
  \end{align}
  Then evaluation of the same weighted binary tree on the two families gives
  \begin{align}
    M_T(A)J&=JM_T(B).
  \end{align}
  Repeated leaf labels and parameters $p=0$ or $p=1$ are allowed.
  No isometry assumption on $J$ is required.
\end{theorem}

\begin{proof}
  \leanok
  \uses{thm:mean_tree_properties, thm:geometric_mean_value_intertwine}
  Induct on the tree. At a leaf labelled $i$, the assertion is the
  assumed identity $A_iJ=JB_i$. At an internal vertex, the induction
  hypotheses give the two intertwining identities for the subtree roots.
  Those roots are positive definite by
  Theorem~\ref{thm:mean_tree_properties}. Applying
  Theorem~\ref{thm:geometric_mean_value_intertwine} at the vertex parameter
  therefore yields
  \begin{align}
    F_p\bigl(M_l(A),M_r(A)\bigr)J
      &=JF_p\bigl(M_l(B),M_r(B)\bigr).
  \end{align}
  By the recursive definition of the root, this is the desired identity
  at that vertex.
\end{proof}
